\section{Overview of the Processing Pipeline} \label{sec:pipeline-overview}

After dataset cleaning, the analysis pipeline begins with a comprehensive preprocessing stage.
This includes timestamp localization to ensure all \gls{eeg} recordings are temporally aligned, accounting for time zone differences and \gls{dst} changes.
Artifacts in the \gls{eeg} signals, such as extreme amplitude events or technical noise, are systematically detected and removed to prevent contamination of downstream analyses.
The continuous \gls{eeg} data is then segmented into short, fixed-length windows (\textit{segments}) and longer, non-overlapping windows (\textit{clips}), with careful labeling of intervals relative to seizure events (e.g., preictal, interictal).
The dataset is partitioned into training and test sets, ensuring sufficient representation of both seizure and non-seizure periods for robust model development and evaluation.

Feature extraction is performed on each segment, yielding a set of quantitative descriptors that capture both temporal and spectral properties of the \gls{eeg}.
These include signal variance, \gls{acfw}, Pearson correlation between channels, and absolute band powers across standard frequency bands.
Some features are used as direct inputs to machine learning models, while all are retained for subsequent cycle analysis.
The extraction process is designed to maximize the interpretability and physiological relevance of the features, supporting both predictive modeling and the investigation of cyclical patterns in seizure risk.

Two main model architectures are employed: a \gls{cnn} that operates on raw \gls{eeg} segments, and an ensemble of \glspl{fb-mlp} that utilize the extracted features.
The \gls{cnn} architecture is adapted to the specific characteristics of the dataset, including its sampling rate and channel configuration, and is regularized to prevent overfitting.
The \gls{fb-mlp} ensemble leverages z-score normalized features and aggregates predictions across multiple independently initialized models.
Both models are trained using only preictal and interictal segments, with careful handling of class imbalance and standardized training procedures to ensure comparability.

Model evaluation is conducted using both standard machine learning metrics and \glspl{ebm} tailored to the clinical context of seizure prediction.
The latter include measures like \gls{tifw} and \gls{eb} sensitivity, which reflect the real-world utility of the models in predicting seizures while minimizing unnecessary alarms.
Thresholds for issuing warnings are optimized to balance \gls{eb} sensitivity and specificity, and all evaluations are performed separately on training and test sets to assess generalization.

To further understand the temporal dynamics of seizure risk, cycle extraction is performed on the feature time series.
This involves gap filling, outlier handling, and bandpass filtering to isolate multidien (multi-day) oscillatory components.
The relationship between cyclical features and event timing is quantified using circular statistics, including \gls{plv} and the Rayleigh test for phase preference.
A non-parametric permutation test is used to compare the phase distributions of seizures and model-predicted \glspl{fp}, with multiple testing correction applied to control the \gls{fdr}.

Finally, the propensity evaluation framework integrates these analyses to determine whether a model’s \glspl{fp} are meaningfully aligned with periods of increased seizure risk.
Qualification criteria are defined for each patient, feature, and model, requiring significant phase locking and predictive performance.
Only when these criteria are met are the phase relationships between \glspl{fp} and seizures interpreted as evidence of the model capturing cyclical seizure propensity, supporting the overall objective of understanding and predicting seizure risk in a clinically relevant manner.
